\section{Introduction}
\label{sec:intro}

Texture partitioning is not just semantic segmentation with unusual labels. The target regions are often defined by repeated local structure, appearance statistics, and transition behavior rather than by named object categories. RWTD and STLD make that distinction explicit, and they also expose the practical bottleneck: texture boundaries can be weak, disconnected, or globally ambiguous even when local evidence is strong~\cite{khan2015stld,khan2017c2f,khan2018learnedstld,mikes2022texturebenchmark}. A system built primarily for broad semantic segmentation is therefore not automatically expected to behave as a texture segmenter.

That is exactly the puzzle raised by SAM-style foundation models. SAM and SAM2 were not trained as texture segmenters~\cite{kirillov2023sam,ravi2024sam2}, yet recent adaptation work such as TextureSAM shows that much stronger texture behavior can be induced by fine-tuning on the RWTD/STLD benchmark family~\cite{cohen2025texturesam,cohen2026texturesam}. This creates a more interesting question than ``can fine-tuning help?'' Fine-tuning working at all suggests that useful texture-relevant structure may already exist inside the frozen system. The scientific question is therefore whether frozen SAM already contains texture understanding, and if so, where that understanding lives and how much of it can be extracted without retraining the backbone or proposal generator.

We study that question through the lens of \emph{recoverability}. A frozen SAM-style system exposes at least two possible loci of texture-relevant evidence. One route explores feature space: coarse frozen features may already organize the target partition before any learned decoder or proposal reasoning. The other explores proposal masks: the generated mask bank may already contain the right fragments, even when the default output fails to assemble them into a coherent texture partition. These routes are complementary rather than hierarchical. The feature route asks what the frozen representation reveals; the proposal-/mask-space route asks what the generated masks reveal. Taken together, they separate latent texture organization from recoverable texture partitions.

\begin{figure*}[t]
\centering
\setlength{\tabcolsep}{8pt}
\renewcommand{\arraystretch}{1.08}
\begin{tabularx}{\textwidth}{@{}>{\raggedright\arraybackslash}m{0.34\textwidth}>{\raggedright\arraybackslash}m{0.62\textwidth}@{}}
\fbox{%
\begin{minipage}[t][4.3cm][t]{0.31\textwidth}
\textbf{Feature-space recovery} \\
\textbf{Frozen evidence:} multiscale SAM \texttt{backbone\_fpn} features. \\
\textbf{Question:} Is texture partition structure already latent in the representation? \\
\textbf{Readout:} pooled coarsest clustering, flip-averaging, positionality diagnostics. \\
\textbf{Role in paper:} diagnostic probe only.
\end{minipage}}
&
\fbox{%
\begin{minipage}[t][4.3cm][t]{0.58\textwidth}
\textbf{Proposal-space recovery} \\
\textbf{Frozen evidence:} SAM-style mask proposals. \\
\textbf{Question:} If the bank already contains the right fragments, can a learned layer recover a coherent partition without retraining the proposal generator? \\
\textbf{Readout:} compatibility reasoning, conservative core selection, and selective repair above the frozen bank. \\
\textbf{Role in paper:} main operational contribution instantiated by \sys{}.
\end{minipage}}
\end{tabularx}
\caption{Two recoverability loci in frozen SAM-style systems. The feature route asks whether texture structure is already latent in frozen features and is used only as a diagnostic probe. The proposal route is the main operational question of the paper: how much task value can be recovered by learning commitment above a frozen proposal bank. Both routes keep the underlying segmenter frozen.}
\label{fig:recoverability_loci}
\end{figure*}


\paragraph{Research questions.} In this paper we ask four questions.
\begin{enumerate}[leftmargin=1.4em]
\item \textbf{RQ1:} Does frozen SAM contain recoverable texture understanding at all?
\item \textbf{RQ2:} What does the frozen feature space reveal about that texture understanding?
\item \textbf{RQ3:} What do the generated masks / proposal masks reveal about that texture understanding?
\item \textbf{RQ4:} Under what regimes do the two routes differ in what they reveal?
\end{enumerate}

\paragraph{Answers and contributions.} The paper's contributions are best read as answers to those questions.
\begin{itemize}[leftmargin=1.2em]
\item To answer \textbf{RQ1} and \textbf{RQ2}, we introduce a feature-space exploration over frozen multiscale features. A lightweight clustering probe, qualitative gallery, and scale / positionality diagnostics show that coarse frozen features already organize nontrivial texture structure.
\item To answer \textbf{RQ1} and \textbf{RQ3}, we introduce a proposal-/mask-space exploration over the generated mask bank. A lightweight learned readout above frozen proposals recovers strong texture partitions under matched RWTD and STLD evaluation without retraining the backbone or proposal generator, improving RWTD common-253 coherence from 0.6163 to 0.7013 ARI against a public TextureSAM rerun while keeping mIoU comparable, and improving STLD common-182 from 0.5140/0.7526 to 0.7195/0.7791.
\item To answer \textbf{RQ4}, we compare what the two routes reveal jointly. Selector controls, oracle decompositions, and the RWTD--STLD contrast show when the generated mask bank already contains a coherent singleton answer and when fragmented mask evidence still requires structured commitment.
\item Across these questions, the paper keeps its claim narrow. We do not argue that frozen SAM fully solves texture segmentation. We argue that frozen SAM contains texture-relevant structure in both frozen features and generated masks, and that the two routes reveal complementary aspects of that structure.
\end{itemize}

\paragraph{Roadmap.} Section~\ref{sec:problem} formalizes recoverability and defines the two loci of frozen evidence. Section~\ref{sec:feature} addresses \textbf{RQ2} through feature-space exploration. Section~\ref{sec:proposal} addresses \textbf{RQ3} through proposal-/mask-space exploration. Section~\ref{sec:results} brings the two routes together to answer \textbf{RQ1} and \textbf{RQ4} through matched RWTD and STLD evidence, learned-selector controls, oracle analysis, and residual failure diagnostics. Section~\ref{sec:discussion} closes the loop on what the paper now knows about frozen SAM's texture understanding and what remains unresolved.
